% GENERATED FILE. Edit canonical lessons, then run npm run generate:books.
% canonical-source-hash: fnv1a64:17eba6c84d841d7b
% canonical-lessons: FR-C319-reveillon, FR-C319-reveillonner, FR-C319-sapin, FR-C319-buche, FR-C319-galette

\chapter{Le réveillon}
\label{ch:fr-le-reveillon}
\hlchaptermodality{319}

\begin{chapteropening}
\textbf{By the end of this chapter:} \emph{I can describe the winter holidays: the late-night feast, celebrating it, the Christmas tree, the yule log and the king cake.}

Complete the last lesson of chapter 319: I can describe the winter holidays: the late-night feast, celebrating it, the Christmas tree, the yule log and the king cake.
\end{chapteropening}

\section[le réveillon]{le réveillon — the festive late-night dinner (24 or 31 December)}

\label{lesson:FR-C319-reveillon}



\begin{quote}
Before the new one: say the French for the garland, then the French for to decorate.
\end{quote}

\subsection*{You'll want to know: le réveillon}
\textbf{le réveillon} — \textquotedblleft{}the festive late-night dinner (24 or 31 December)\textquotedblright{}.

\begin{quote}
Toute la famille passe le réveillon chez ma tante.
\end{quote}

\begin{quote}
\emph{The whole family spends the festive dinner at my aunt's.}
\end{quote}

It falls on the evenings of 24 and 31 December.

\begin{cousinweb}
\emph{réveillon} is built on \emph{réveiller} \textquotedblleft{}to wake up\textquotedblright{}, from \emph{veiller} \textquotedblleft{}to stay awake\textquotedblright{}, which is inherited from Latin \emph{vigilāre} \textquotedblleft{}to keep watch\textquotedblright{} (Lewis \& Short).
\end{cousinweb}

\begin{grammarlens}[title={What you've built}]
One more word for telling, explaining and describing at B1.
\end{grammarlens}

\subsection*{Your turn}
\begin{itemize}
\raggedright
  \item \emph{Say it:} \emph{le réveillon}
  \item \emph{Say it:} \emph{le réveillon}, once more
  \item \emph{Say it:} a sentence of your own with \emph{le réveillon}
  \item \emph{Recall:} say the French for the garland, then the French for to decorate, then say \emph{le réveillon} again
\end{itemize}

\begin{tcolorbox}[breakable,colback=teal!4,colframe=teal!35!black,title={Before you move on}]
What does le réveillon mean? (\textquotedblleft{}The festive late-night dinner (24 or 31 December)\textquotedblright{}.) Say it once more.
\end{tcolorbox}
\section[réveillonner]{réveillonner — to have the festive late-night dinner}

\label{lesson:FR-C319-reveillonner}



\begin{quote}
Before the new one: say the French for to decorate, then the French for the festive late-night dinner.
\end{quote}

\subsection*{You'll want to know: réveillonner}
\textbf{réveillonner} — \textquotedblleft{}to have the festive late-night dinner\textquotedblright{}.

\begin{quote}
Nous avons réveillonné jusqu'à trois heures du matin.
\end{quote}

\begin{quote}
\emph{We had our festive dinner until three in the morning.}
\end{quote}

On 24 December or 31 December.

\begin{cousinweb}
\emph{réveillonner} is built on \emph{réveillon}, from \emph{réveiller} \textquotedblleft{}to wake up\textquotedblright{}, which goes back to Latin \emph{vigilāre} \textquotedblleft{}to keep watch\textquotedblright{} (Lewis \& Short).
\end{cousinweb}

\begin{grammarlens}[title={What you've built}]
One more word for telling, explaining and describing at B1.
\end{grammarlens}

\subsection*{Your turn}
\begin{itemize}
\raggedright
  \item \emph{Say it:} \emph{réveillonner}
  \item \emph{Say it:} \emph{réveillonner}, once more
  \item \emph{Say it:} a sentence of your own with \emph{réveillonner}
  \item \emph{Recall:} say the French for to decorate, then the French for the festive late-night dinner, then say \emph{réveillonner} again
\end{itemize}

\begin{tcolorbox}[breakable,colback=teal!4,colframe=teal!35!black,title={Before you move on}]
What does réveillonner mean? (\textquotedblleft{}To have the festive late-night dinner\textquotedblright{}.) Say it once more.
\end{tcolorbox}
\section[le sapin]{le sapin — the fir tree; the Christmas tree}

\label{lesson:FR-C319-sapin}



\begin{quote}
Before the new one: say the French for the festive late-night dinner, then the French for to have the festive late-night dinner.
\end{quote}

\subsection*{You'll want to know: le sapin}
\textbf{le sapin} — \textquotedblleft{}the fir tree; the Christmas tree\textquotedblright{}.

\begin{quote}
Les enfants décorent le sapin.
\end{quote}

\begin{quote}
\emph{The children decorate the Christmas tree.}
\end{quote}

Decorated in December, often with \emph{des guirlandes}.

\begin{cousinweb}
\emph{sapin} is inherited from Latin \emph{sappīnus} \textquotedblleft{}a fir tree\textquotedblright{} (EtymDB), probably of Gaulish origin.
\end{cousinweb}

\begin{grammarlens}[title={What you've built}]
One more word for telling, explaining and describing at B1.
\end{grammarlens}

\subsection*{Your turn}
\begin{itemize}
\raggedright
  \item \emph{Say it:} \emph{le sapin}
  \item \emph{Say it:} \emph{le sapin}, once more
  \item \emph{Say it:} a sentence of your own with \emph{le sapin}
  \item \emph{Recall:} say the French for the festive late-night dinner, then the French for to have the festive late-night dinner, then say \emph{le sapin} again
\end{itemize}

\begin{tcolorbox}[breakable,colback=teal!4,colframe=teal!35!black,title={Before you move on}]
What does le sapin mean? (\textquotedblleft{}The fir tree; the Christmas tree\textquotedblright{}.) Say it once more.
\end{tcolorbox}
\section[la bûche]{la bûche — the log; the yule log cake}

\label{lesson:FR-C319-buche}



\begin{quote}
Before the new one: say the French for to have the festive late-night dinner, then the French for the fir tree.
\end{quote}

\subsection*{You'll want to know: la bûche}
\textbf{la bûche} — \textquotedblleft{}the log; the yule log cake\textquotedblright{}.

\begin{quote}
Pour le dessert, nous avons mangé une bûche au chocolat.
\end{quote}

\begin{quote}
\emph{For dessert, we ate a chocolate yule log.}
\end{quote}

A rolled cake shaped like a log, eaten in December.

\begin{cousinweb}
\emph{bûche} comes from Old French \emph{busche}, inherited from Vulgar Latin \emph{busca} \textquotedblleft{}firewood\textquotedblright{} (EtymDB), usually traced to a Germanic word.
\end{cousinweb}

\begin{grammarlens}[title={What you've built}]
One more word for telling, explaining and describing at B1.
\end{grammarlens}

\subsection*{Your turn}
\begin{itemize}
\raggedright
  \item \emph{Say it:} \emph{la bûche}
  \item \emph{Say it:} \emph{la bûche}, once more
  \item \emph{Say it:} a sentence of your own with \emph{la bûche}
  \item \emph{Recall:} say the French for to have the festive late-night dinner, then the French for the fir tree, then say \emph{la bûche} again
\end{itemize}

\begin{tcolorbox}[breakable,colback=teal!4,colframe=teal!35!black,title={Before you move on}]
What does la bûche mean? (\textquotedblleft{}The log; the yule log cake\textquotedblright{}.) Say it once more.
\end{tcolorbox}
\section[la galette]{la galette — the king cake; a flat round cake}

\label{lesson:FR-C319-galette}



\begin{quote}
Before the new one: say the French for the fir tree, then the French for the log.
\end{quote}

\subsection*{You'll want to know: la galette}
\textbf{la galette} — \textquotedblleft{}the king cake; a flat round cake\textquotedblright{}.

\begin{quote}
En janvier, on mange la galette des rois.
\end{quote}

\begin{quote}
\emph{In January, people eat the king cake.}
\end{quote}

\emph{La galette des rois} hides a small figure; whoever finds it wears the paper crown.

\begin{cousinweb}
\emph{galette} is built on \emph{galet} \textquotedblleft{}a flat pebble\textquotedblright{}, a diminutive of an old word that is usually traced to Gaulish (EtymDB lists no parent).
\end{cousinweb}

\begin{grammarlens}[title={What you've built}]
One more word for telling, explaining and describing at B1.
\end{grammarlens}

\subsection*{Your turn}
\begin{itemize}
\raggedright
  \item \emph{Say it:} \emph{la galette}
  \item \emph{Say it:} \emph{la galette}, once more
  \item \emph{Say it:} a sentence of your own with \emph{la galette}
  \item \emph{Recall:} say the French for the fir tree, then the French for the log, then say \emph{la galette} again
\end{itemize}

\begin{tcolorbox}[breakable,colback=teal!4,colframe=teal!35!black,title={Before you move on}]
What does la galette mean? (\textquotedblleft{}The king cake; a flat round cake\textquotedblright{}.) Say it once more.
\end{tcolorbox}
